% !TeX root = ../../Book3.tex
% 纲要标题：### 20.1 R_k 与 S_k 条件
% 纲要位置：计划纲要.md，第 2116 行。

\section{\texorpdfstring{\(R_k\)}{Rₖ} 与 \texorpdfstring{\(S_k\)}{Sₖ} 条件}
\label{b3:sec:2001}

正规性的局部检验同时涉及低余维处的环结构与较高余维处的深度。先分别固定 \(R_k\) 和 \(S_k\) 的量词，再用算例说明二者不能互相替代。


% 纲要内容：
%
% * 余维上的正则性
% * 深度条件的局部化表述
% * 余维一和高余维承担的不同任务
%

\subsection{余维上的正则性}
\label{b3:subsec:2001:01}

整闭性表面上要求检查分式域中的所有首一方程，正则性则比较局部环的维数与极大理想生成元数。Serre 判据把前一种条件分成两项：低余维处检查正则性，较高余维处检查深度。先把两项条件的量词固定下来。

\begin{definition}\label{b3:def:serre-rk}
设 \(R\) 为交换 Noether 环，\(k\ge0\) 为整数。若对每个满足 \(\operatorname{ht}\mathfrak p\le k\) 的素理想 \(\mathfrak p\)，局部环 \(R_{\mathfrak p}\) 都正则，则称 \(R\) 满足 \term[Serre-R-tiaojian]{Serre 条件 \(R_k\)}[Serre's condition \(R_k\)]。
\end{definition}

这里的高度是 \(\dim R_{\mathfrak p}\)，不是 \(\dim R/\mathfrak p\)。在仿射整簇中，它衡量所对应子簇的余维。条件 \(R_0\) 表示极小素点处的局部环为域；条件 \(R_1\) 还要求每个高度一素点处的局部环为 DVR，因为一维正则局部环的极大理想由一个元素生成。

\ProseParagraphBreak
条件 \(R_0\) 只保证在各分支的一般位置没有幂零元，不能排除集中在特殊点的幂零元。例如
\(k[x,y]/(x^2,xy)\) 在唯一极小素理想 \((x)\) 处把 \(y\) 变成单位，从而 \(x=0\)，局部化为域 \(k(y)\)；但原环中的 \(x\) 非零且平方为零。

\subsection{深度条件的局部化表述}
\label{b3:subsec:2001:02}

回顾\secref{b3:sec:1604}：交换 Noether 环 \(R\) 满足 \(S_k\)，是指对每个素理想 \(\mathfrak p\)，都有
\[
\operatorname{depth}R_{\mathfrak p}\ge\min\{k,\dim R_{\mathfrak p}\}.
\]
特别地，\(S_1\) 排除嵌入关联素理想；\(S_2\) 在所有维数至少二的局部环中要求能连续取两个非零因子。在维数一处，它只要求一个；维数零处没有额外要求。

\begin{proposition}\label{b3:prop:serre-conditions-localize}
设 \(R\) 为交换 Noether 环，\(k\ge0\) 为整数。条件 \(R_k\)、\(S_k\) 均在局部化后保持；它们也都可在各个极大理想处的局部环上检查。若 \(R\) 为 CM 环，则它满足每个 \(S_k\)。
\end{proposition}
\begin{proof}
若 \(S^{-1}\mathfrak p\) 是 \(S^{-1}R\) 的素理想，则
\((S^{-1}R)_{S^{-1}\mathfrak p}\cong R_{\mathfrak p}\)；其以下的素理想链也都避开 \(S\)，所以高度不变。于是局部条件完全相同。每个素理想包含在某个极大理想中，进一步局部化也得到同一个 \(R_{\mathfrak p}\)，证明检查极大理想处足够。CM 环在各处深度等于维数，因此满足所列全部不等式。
\end{proof}

\begin{proposition}\label{b3:prop:reduced-serre-criterion}
对交换 Noether 环 \(R\)，约化性等价于 \(R_0\) 与 \(S_1\) 同时成立。
\end{proposition}
\begin{proof}
若 \(R\) 约化，则其极小素点处为约化零维局部环，因而为域。
\cref{b3:prop:reduced-ring-no-embedded-primes}和\cref{b3:thm:serre-s1-associated}给出 \(S_1\)。

反之，若幂零根 \(N\) 非零，则它作为有限 \(R\) 模有一个关联素理想 \(\mathfrak p\)。因为 \(N\subseteq R\)，有 \(\mathfrak p\in\operatorname{Ass}R\)；\(S_1\) 使它为极小素理想。于是 \(N_{\mathfrak p}\ne0\)，但 \(R_0\) 使 \(R_{\mathfrak p}\) 为域，其幂零根只能为零，矛盾。
\end{proof}

\subsection{余维一与高余维的条件比较}
\label{b3:subsec:2001:03}

\(R_1\) 与 \(S_2\) 不能互相替代。一维尖点局部环
\(k[t^2,t^3]_{(t^2,t^3)}\) 是 CM 环，故满足 \(S_2\)；它却不是 DVR，所以不满足 \(R_1\)。另一类例子由两张平面只在一个点相交得到。

\begin{example}\label{b3:exm:r1-not-s2}
设 \(k\) 为域，
\[
T=k[x,y,z,w]_{(x,y,z,w)},\qquad
R=T/\bigl((x,y)\cap(z,w)\bigr).
\]
则 \(R\) 为二维约化局部环，满足 \(R_1\)，却有 \(\operatorname{depth}R=1\)，因而不满足 \(S_2\)。
\end{example}
\begin{proof}
两个极小分支的商均为二维正则局部环，且它们的和为 \(T\) 的极大理想，故两分支只在闭点相遇，\(\dim R=2\)。设闭点为 \(\mathfrak m\)。若 \(\mathfrak p\ne\mathfrak m\)，至少有一个坐标在 \(\mathfrak p\) 外；把它逆转后，相对的两个坐标由交理想中的乘积关系变为零，局部环成为其中一个正则分支的局部化。
\cref{b3:thm:regular-local-localization}使它正则，故 \(R_1\) 成立。

自然差映射给出正合列
\[
0\longrightarrow R\longrightarrow T/(x,y)\oplus T/(z,w)
\xrightarrow{(a,b)\mapsto a-b}k\longrightarrow0.
\]
中间模上 \(x+z,y+w\) 依次为非零因子，在两个直和分量中分别是两个剩余变量，所以它的深度为二；右端的深度为零。\cref{b3:thm:depth-lemma}给出 \(\operatorname{depth}R=1\)：深度至少一来自中间模深度，若至少二，则对上述正合列应用 \(\operatorname{Hom}_R(k,-)\) 的长正合列会使从非零的 \(\operatorname{Hom}_R(k,k)\) 到零的 \(\operatorname{Ext}_R^1(k,R)\) 为单射，矛盾。
\end{proof}

余维一检查没有看见只发生在闭点的相交；深度二要求恰好排除了这个障碍。Serre 判据将证明，对于整闭性，这两类检查合起来已经足够。
